博弈题先列行动顺序、可见信息、支付与平局规则，再写单方偏离后的真实收益。私人信息或成交筛选会改变条件分布；优化报价前须先确定面对的是哪一类对象。个人均衡与总成本最优应分别计算。

零和混合策略可通过双方的收益保证证明最优；支持内等收益只产生候选，还要检查所有支持外行动。有限取物游戏可用强归纳分类局面，或设计每次都合法的对称回应，不能仅凭几个小例子断言周期。

构造题也要检验实施限制：匹配需核对邻集条件，同余组需检查相容性，称量编码除了区分状态，还必须满足每次两盘的平衡要求。
\par\smallskip
